%Version 3.1 December 2024
% See section 11 of the User Manual for version history
%
%%%%%%%%%%%%%%%%%%%%%%%%%%%%%%%%%%%%%%%%%%%%%%%%%%%%%%%%%%%%%%%%%%%%%%
%%                                                                 %%
%% Please do not use \input{...} to include other tex files.       %%
%% Submit your LaTeX manuscript as one .tex document.              %%
%%                                                                 %%
%% All additional figures and files should be attached             %%
%% separately and not embedded in the \TeX\ document itself.       %%
%%                                                                 %%
%%%%%%%%%%%%%%%%%%%%%%%%%%%%%%%%%%%%%%%%%%%%%%%%%%%%%%%%%%%%%%%%%%%%%

\documentclass[pdflatex,sn-mathphys-num]{sn-jnl}% Math and Physical Sciences Numbered Reference Style

%%%% Standard Packages
\usepackage{graphicx}%
\usepackage{multirow}%
\usepackage{amsmath,amssymb,amsfonts}%
\usepackage{amsthm}%
\usepackage[title]{appendix}%
\usepackage{xcolor}%
\usepackage{textcomp}%
\usepackage{manyfoot}%
\usepackage{booktabs}%
\usepackage{algorithm}%
\usepackage{algorithmicx}%
\usepackage{algpseudocode}%
\usepackage{listings}%
\usepackage{array}%
%% ragged right paragraph column for the audit tables
\newcolumntype{L}[1]{>{\raggedright\arraybackslash}p{#1}}
%% mathrsfs is deliberately not loaded: it is unused here, and loading it makes
%% every small math size request the rsfs font.

%% ---------------------------------------------------------------------
%% ADDED FOR COMPILATION ONLY -- delete these six lines once figs/*.png are
%% in place. The three figures are not in this repository; without this the
%% document fails to build on a machine that does not have them. With it,
%% a missing figure renders as a labelled placeholder box and the rest of
%% the document still compiles. Nothing else in the file depends on it.
\newcommand{\figorbox}[2]{\IfFileExists{#2}{\includegraphics[width=#1]{#2}}%
  {\fbox{\parbox[c][45mm][c]{#1}{\centering\ttfamily\small MISSING FIGURE\\#2}}}}
%% ---------------------------------------------------------------------

\theoremstyle{thmstyleone}%
\newtheorem{theorem}{Theorem}%
\newtheorem{proposition}[theorem]{Proposition}%

\theoremstyle{thmstyletwo}%
\newtheorem{example}{Example}%
\newtheorem{remark}{Remark}%

\theoremstyle{thmstylethree}%
\newtheorem{definition}{Definition}%

\raggedbottom

%% The class does \LoadClass[twoside,fleqn]{article}, so odd and even pages get
%% mirrored margins. Passing oneside does not help, because the class applies its
%% own twoside afterwards. Average the two instead.
\AtBeginDocument{%
  \addtolength{\oddsidemargin}{\evensidemargin}%
  \setlength{\oddsidemargin}{0.5\oddsidemargin}%
  \setlength{\evensidemargin}{\oddsidemargin}%
}

%% \bmhead is a level-5 heading while the class sets bookmarksdepth=5, so the
%% back matter jumps four bookmark levels at once and hyperref complains.
\hypersetup{bookmarksdepth=3}

%% DOIs and URLs cannot hyphenate, so the surrounding bibliography lines are set
%% loose. \sloppy permits the loose line but leaves \hbadness at 1000; both are
%% needed, scoped to the bibliography and applied after natbib's own hooks.
\AtBeginDocument{%
  \let\snOrigThebibliography\thebibliography
  \renewcommand{\thebibliography}[1]{%
    \snOrigThebibliography{#1}\sloppy\hbadness=10000\relax}%
}

\begin{document}

\title[Causal Evaluation of Decomposition Ensemble Stock Index Forecasting]{Causal Evaluation of Decomposition Ensemble Stock Index Forecasting: Reported Accuracy Depends on the Decomposition Seeing the Future}

%% AUTHORS: copied from the original Group 17 submission as a placeholder.
%% Confirm the list and order, add the members who carried out this study,
%% and supply every e-mail address (the \email lines marked TODO) before
%% submission. No address below has been invented

\author[1]{\fnm{Prayag} \sur{Sharma}}
\email{prayag.sharma2024@vitstudent.ac.in}

\author[1]{\fnm{Veer Vikram} \sur{Saxena}}
\email{veer.vikramsaxena2024@vitstudent.ac.in}

\author[1]{\fnm{Ayush} \sur{Srivastava}}
\email{ayush.srivastava2024@vitstudent.ac.in}

\author*[2]{\fnm{Ravi} \sur{Tiwari}}
\email{ravi.tiwari@vit.ac.in}

\affil[1]{%
\orgdiv{School of Computer Science and Engineering (SCOPE)},
\orgname{Vellore Institute of Technology (VIT)},
\orgaddress{\city{Chennai}, \state{Tamil Nadu}, \country{India}}%
}

\affil[2]{%
\orgdiv{School of Electronics Engineering (SENSE)},
\orgname{Vellore Institute of Technology (VIT)},
\orgaddress{\city{Chennai}, \state{Tamil Nadu}, \country{India}}%
}

\abstract{Decomposition ensemble models, which split a price series into
components with empirical or variational mode decomposition and forecast each
component with a neural network, report daily index forecasts with percentage
errors well below one tenth of one percent. A decomposition is not causal,
however: the value of a component on a given day depends on the days that
follow it, so a series decomposed in full before the training and test split
carries the test prices into the test inputs. We measure how much of the
reported accuracy depends on this. Reproducing a published model that combines
improved complete ensemble empirical mode decomposition with adaptive noise, a
second variational mode decomposition stage tuned by particle swarm
optimisation, and a bidirectional long short term memory network with self
attention feeding either a temporal convolutional head or a tiny recursive
head, we compare full series decomposition with a walk forward protocol that
redecomposes the available history at every forecast origin, on the daily
highest and lowest prices of the S\&P~500 and the Shanghai Stock Exchange
Composite Index over ten years. Before any network is trained, a linear
regression on the full series inputs explains 51 to 60 percent of the next
day's price change, and none of it under walk forward decomposition. Trained
on identical code, seeds and budgets, the networks reach a mean absolute
percentage error of 0.07 to 0.13 percent and a coefficient of determination
above 0.999 under full series decomposition and beat persistence on every
series by the modified Diebold and Mariano test; under walk forward
decomposition their error is 4.9 to 10.4 times larger and they fall
significantly behind persistence on every series. The reported advantage of
the convolutional head over the recursive head holds only under full series
decomposition. The evidence covers two indices, a one day horizon and one seed
per network. In reproducing the pipeline we identified and corrected twelve
implementation issues, which we document, and we propose two inexpensive
checks, an input diagnostic and a leakage ratio, that make the dependence of a
headline result on its protocol visible.}

\keywords{Deep learning, Financial time series, Look ahead bias, Boundary issue, Walk forward validation, Diebold and Mariano test}

\maketitle

\section{Introduction}\label{sec:intro}

Forecasting daily stock index levels is a problem where small reported gains
attract a great deal of attention, because a model that is even slightly more
accurate than its competitors has obvious value to anyone managing risk.
Machine learning for stock market prediction is by now a large literature, with
well over a hundred journal articles reviewed in a single survey~\cite{kumbure2022},
and over the past decade one design has come to dominate the applied part of
it: split the price series into simpler subseries with a signal decomposition,
forecast each subseries with its own neural network, and add the forecasts back
together. Empirical mode decomposition (EMD) and its interpolation and noise
assisted variants~\cite{akimaemd2023,gruceemdan2023,ceemdanstock2025} and
variational mode decomposition (VMD)~\cite{jiang2025,wangwang2026} are the usual
front ends, and recurrent and convolutional networks, often with attention, are
the usual back ends~\cite{gong2024,zhao2025,ning2025}. Reported accuracies are
often high, with coefficients of determination close to
one~\cite{ning2025,jiang2025}.

Numbers of that kind deserve a second look for a specific reason. A signal
decomposition is not a causal filter: the value of a component at time $t$
depends on the whole stretch of data that was decomposed, including samples
after $t$. When a study decomposes the complete series first and only then
separates training from test data, every input the model sees at test time has
been computed with knowledge of the prices it is asked to predict. Leakage of
this general kind has produced overoptimistic conclusions in many fields that
adopted machine learning~\cite{kapoor2023}. For decomposition based
forecasters it is well documented in wind power forecasting, where it is known
as the boundary issue~\cite{chen2022}; it has been shown to create an illusion
of accuracy in water quality prediction~\cite{emdleak2024}; and several recent
studies design decomposition schemes that keep future samples out of the model
inputs~\cite{wang2024,yao2025,pang2026}. Look ahead bias is also a recognised
concern in financial backtesting, most recently for signals extracted by large
language models~\cite{glasserman2024}. Yet this particular route for future
information has received far less systematic attention in equity forecasting,
where the same pipelines are used and where, under the random walk and
efficient market hypotheses, future prices cannot be systematically
predicted~\cite{pernagallo2025}, even though machine learning classifiers have
been reported to beat random choice on longer horizon relative
returns~\cite{finrel2024}.

This paper measures how much of the apparent skill of a decomposition ensemble
stock forecaster survives when the decomposition is made strictly causal. We
work with a published redecomposition model that combines improved complete
ensemble EMD with adaptive noise (ICEEMDAN), a second VMD stage tuned by
particle swarm optimisation (PSO), and a bidirectional long short term memory
(BiLSTM) network with self attention feeding a temporal convolutional network
(TCN)~\cite{gong2024}, together with a variant that replaces the convolutional
head with a Tiny Recursive Model (TRM)~\cite{jolicoeur2025}. On the daily
highest and lowest prices of the S\&P~500 and the Shanghai Stock Exchange
Composite Index (SSEC) we compare two protocols that differ in one respect
only: whether the series is decomposed once in full or redecomposed from the
available history at every forecast origin. The contributions are as follows.
\begin{enumerate}
\item We show, before any network is trained, that full series decomposition
places information about the next day's price change into the model inputs,
and that walk forward decomposition removes it.
\item We quantify the accuracy attributable to decomposing the full series,
as the ratio of out of sample error under walk forward decomposition to the
error under full series decomposition, measured on identical test dates for
every index, price series and network.
\item We evaluate both forecasters against persistence, drift, moving average
and exponential smoothing benchmarks with the modified Diebold and Mariano
(MDM) test~\cite{hewamalage2022,mdmindex2023}, so that any claimed advantage is
judged against forecasts that require no model at all.
\item In reproducing the pipeline we identified fourteen issues. Twelve were
corrected, several of which would on their own have changed the central
comparison, and two are modelling choices of the source model that we keep
and report. Section~\ref{sec:audit} documents each one, so that the
reproduction can be checked rather than trusted.
\item We give both network heads identical training budgets, schedules and
postprocessing, which the source comparison did not, and report whether the
ranking of the two heads depends on the protocol.
\end{enumerate}
Section~\ref{sec:related} reviews the related work, Section~\ref{sec:data}
describes the data and the two protocols, Section~\ref{sec:setup} the
networks, training, measures and the audit of the reproduced pipeline,
Section~\ref{sec:results} the results, Section~\ref{sec:discussion} their
interpretation and limits, and Section~\ref{sec:conclusion} concludes.

\section{Related work and its limitations}\label{sec:related}

This section covers three topics in turn: the decomposition ensemble design
and its recent extensions, the evidence that a decomposition can carry future
information into a forecaster, and the benchmarks and network heads that this
study uses.

\subsection{Decomposition ensemble forecasting}
EMD represents a nonstationary series as a sum of intrinsic mode functions
extracted by an iterative, spline based sifting procedure whose behaviour at
the boundaries of the data needs particular care~\cite{emdtutorial2023}. Plain
EMD is prone to mode mixing, and ICEEMDAN is one of the variants adopted to
counter it~\cite{iceemdanbridge2022}; such noise assisted methods decompose
noise added copies of the signal, together with the added noise itself, using
classical EMD~\cite{ceemdanpaa2022}. VMD extracts a preset number of modes and
cannot choose its input parameters by itself~\cite{vmdmssr2022}, so they are
frequently tuned by an optimiser such as PSO~\cite{psovmdtcn2023}.
Decomposition based hybrids have become very popular~\cite{chen2022}, and
recent work continues to extend them: EMD with long short term memory (LSTM)
networks for the KSE-100 index~\cite{akimaemd2023}, complete ensemble EMD with
adaptive noise (CEEMDAN) with wavelet denoising and gated recurrent units for
the S\&P~500 and CSI~300~\cite{gruceemdan2023}, CEEMDAN with LSTM and
backpropagation networks for the DAX~\cite{ceemdanstock2025}, a second VMD
stage applied to the high frequency modes of a CEEMDAN decomposition for wind
power~\cite{ning2025}, VMD followed by convolutional, bidirectional recurrent
and attention layers for natural gas futures prices~\cite{jiang2025}, adaptive
VMD with feature selection and a BiLSTM for stock prices~\cite{wangwang2026},
PSO tuned VMD with an attention TCN for electricity load~\cite{psovmdtcn2023},
wavelet denoising combined with an attention BiLSTM, an autoregressive
integrated moving average (ARIMA) model and gradient boosting for the CSI~300
index~\cite{zhao2025}, metaheuristic tuning of decomposition aided attention
recurrent networks for solar and wind generation~\cite{damasevicius2024}, and
decomposition blocks placed inside general forecasting
architectures~\cite{kreuzer2025}.

\subsection{Future information in decomposition}
Chen \emph{et al.} reviewed decomposition based wind power models and noted
that existing studies usually decompose data that would not be available at
forecast time, so the reported accuracy is higher than what can be achieved in
practice~\cite{chen2022}. Yang \emph{et al.} found the same effect in EMD
based water quality prediction, where decomposing before the split let test
information into training and produced an illusion of accuracy, and they
proposed sliding window and repeated decomposition schemes to remove
it~\cite{emdleak2024}. Other remedies include cascade decomposition designed to
avoid information leakage~\cite{wang2024}, rolling decomposition in which the
point to be predicted never enters the decomposition~\cite{yao2025}, and
explicit guidance on preprocessing that prevents future data from leaking into
wind speed forecasts~\cite{pang2026}. The walk forward protocol used here
follows the same principle and matches recommended practice for partitioning
time series in forecast evaluation~\cite{hewamalage2022}. What none of this
work provides is a measurement, on equity indices, of how much of a published
model's accuracy the leak accounts for; that is the gap this paper addresses.

\subsection{Benchmarks and network heads}
Under the random walk and efficient market hypotheses future index levels
cannot be systematically predicted~\cite{pernagallo2025}, so persistence is a
demanding benchmark, and simple models deserve a place in any comparison: a one
layer linear model outperformed sophisticated Transformer based forecasters on
nine real datasets~\cite{zeng2023}, and a large comparison on the M3 competition
series found that deep learning adds value mainly in combination and at a
considerable computational cost~\cite{makridakis2022}. Machine learning
studies often adopt evaluation practices that make uncompetitive methods look
competitive, and we follow published guidance on data partitioning, error
measures and statistical testing~\cite{hewamalage2022}; a systematic survey of
deep learning in the stock market likewise calls for more reproducible,
explainable and accountable models~\cite{olorunnimbe2022}. Forecasts are
compared with the MDM test, which has been used to compare daily index price
forecasts~\cite{mdmindex2023}. The encoder combines a BiLSTM with attention, a
combination covered in recent reviews of recurrent
architectures~\cite{rnnreview2024} and the one used by the source
model~\cite{gong2024}. One head is a TCN, which recent decomposition based
forecasters also pair with attention~\cite{psovmdtcn2023,emdleak2024}; the
other adapts the recursive latent refinement of the TRM~\cite{jolicoeur2025}.

\section{Data and decomposition protocols}\label{sec:data}

The experiment is defined by the price series, the two stage decomposition
applied to them, and the two protocols under which that decomposition is run.
Five subsections describe them: the series, the decomposition, protocol A,
protocol B, and how results under the two protocols are matched.

\subsection{Price series}
The data are the daily highest and lowest prices of the S\&P~500 (ticker
\texttt{\textasciicircum GSPC}) and of the SSEC (\texttt{000001.SS}) over a
ten year window, downloaded once as unadjusted values and cached, so that
every experiment uses the same array. Rows with a missing highest or lowest
price were removed jointly, keeping the two series of an index aligned by
trading day. The S\&P~500 series contain $2513$ observations and the SSEC
series $2425$. Each series is forecast one step ahead from a window of the
previous $30$ values. Of the windowed samples, the first $80\%$ form the
training set and the rest the test set, giving $497$ test days for the
S\&P~500 and $479$ for the SSEC.

\subsection{Two stage decomposition}
The first stage applies ICEEMDAN~\cite{colominas2014} to the price
series, producing intrinsic mode functions ordered from highest to lowest
frequency and a final residual. Sifting stops when the normalised squared
difference between successive iterates falls below $0.2$, with at most $100$
sifting iterations and $10$ modes. The second stage decomposes the highest
frequency mode again with VMD~\cite{dragomiretskiy2014}, with the number of modes $K$
and the bandwidth penalty $\alpha$ chosen by PSO~\cite{psosurvey2022} over
$K\in[3,10]$ and $\alpha\in[100,5000]$ (twenty particles, thirty iterations,
inertia $0.7$, both acceleration coefficients $1.5$). The $K$ VMD modes, the
remaining ICEEMDAN modes and the residual form the $C$ subseries that are
forecast, with $C$ equal to $16$ or $17$ here.

Two properties of this pipeline, inherited from the source model and kept
unchanged so that the reproduction stays faithful, matter for the
interpretation of the results. First, the swarm minimises the reconstruction
error of the decomposed mode, and that error falls monotonically as $K$ grows
and $\alpha$ shrinks. The search therefore returned the bounds $K=10$ and
$\alpha=100$ on all eight decompositions, so the optimisation step does not in
practice select anything. Second, the noise added in each ICEEMDAN realisation
has a fixed standard deviation of $0.2$ price units, between
$2\times10^{-4}$ and $8\times10^{-4}$ of the standard deviation of the
training series. This amplitude is small, although it still alters the
sifting: two runs that differ only in their random seed disagree by up to
$1.5\%$ of the price range in individual components. Our results should
therefore be read as describing ICEEMDAN at this noise level.

\subsection{Protocol A: full series decomposition}
Let $x_1,\dots,x_N$ be a price series, $\mathcal{D}$ the two stage
decomposition, and $c_k(t)$ the value of component $k$ at time $t$. A
forecast of $x_t$ is formed from component forecasts,
\begin{equation}
\hat{x}_t=\sum_{k=1}^{C}\hat{c}_k(t),\qquad
\hat{c}_k(t)=f_k\bigl(c_k(t-30),\dots,c_k(t-1)\bigr),
\label{eq:ensemble}
\end{equation}
where each $f_k$ is a separately trained network. The protocols differ only in
how the component values entering $f_k$ are computed. Under protocol A the
whole series is decomposed once, $c_k(\cdot)=\mathcal{D}_k(x_{1:N})$, and the
resulting components are then divided into training and test windows. Both the
swarm search and every component value, including those inside test windows,
depend on all $N$ observations. This is the protocol of the source model and of
much of the literature reviewed above.

\subsection{Protocol B: walk forward decomposition}
Under protocol B the training portion alone is decomposed, and the swarm
search runs on it alone; the resulting $K$, $\alpha$ and $C$ are then held
fixed. For every test origin $t$ the available history is decomposed afresh
and the last $30$ rows of that decomposition form the input,
\begin{equation}
c_k^{(t)}(s)=\mathcal{D}_k\bigl(x_{a_t:t-1}\bigr)(s),\qquad s\le t-1,
\label{eq:walkforward}
\end{equation}
so no value at or after $t$ is ever used to build an input. VMD needs an input
of even length, and $a_t\in\{1,2\}$ is chosen to provide one: when the history
has odd length its oldest observation is left out, so the newest observation
$x_{t-1}$ is always kept. Decomposing histories of different lengths can yield
a different number of intrinsic mode functions; surplus low frequency columns
are added to the residual and a shortfall is filled with a zero column
immediately before it, which leaves the sum of the components unchanged. Each
series needs one decomposition per test day plus one, $498$ for the S\&P~500
and $480$ for the SSEC.

\subsection{Matching the two protocols}
Two practical differences between the protocols were measured rather than
assumed away. Protocol A was run with $50$ noise realisations, as in the
source model, and protocol B with $20$ to keep the $1956$ walk forward
decompositions tractable. Decomposing the training portion of the highest
price series both ways, the largest difference in any component was $1.3\%$
of the price range, comparable to the $1.5\%$ separating two realisation runs
of $50$ that differ only in their seed, so the smaller ensemble adds about as
much variation as a change of seed and no more. In addition, the full series
decomposition trims the series to an even length at its end, so the two
protocols place their split one day apart. All comparisons between protocols
are therefore computed on the dates both protocols test, $496$ for the
S\&P~500 and $478$ for the SSEC.

\section{Experimental setup}\label{sec:setup}

Five subsections describe the setup: the two networks, their training, the
evaluation measures, the issues identified in the reproduced pipeline and how
each was resolved, and what a complete run of the study consists of.

\subsection{Networks}
Both forecasters share an encoder: a two layer BiLSTM with $64$ units per
direction and dropout of $0.1$ between layers~\cite{rnnreview2024}, followed by
four head self attention with a residual connection and layer normalisation,
as in the source model~\cite{gong2024}. The \emph{TCN head} applies a causal
dilated convolution (kernel $3$, dilation $2$) with a rectified linear unit,
dropout of $0.1$ and a residual linear path, and reads the forecast from the
last position. The \emph{TRM head} projects the last encoder state and refines
a latent state $z$ and an answer state $y$ recursively, applying a two layer
latent network six times per cycle and an answer network once, for three
cycles of which only the last is differentiated~\cite{jolicoeur2025}.

\subsection{Training}
One network is trained per component, so each protocol, index, price series
and head involves $16$ or $17$ networks and the study trains $262$ in total.
Each input window is expressed relative to its own last value, and the network
$g_k$ forecasts the change from that value to the next one,
\begin{equation}
\tilde{u}_j=\frac{c_k(t-31+j)-c_k(t-1)}{\sigma_{\mathrm{in}}},\quad j=1,\dots,30,
\qquad
\hat{c}_k(t)=c_k(t-1)+\sigma_{\mathrm{out}}\,g_k(\tilde{u}),
\label{eq:scaling}
\end{equation}
where $\sigma_{\mathrm{in}}$ and $\sigma_{\mathrm{out}}$ are the standard
deviations of the input deviations and of the one step changes over the
training windows. Component forecasts from~\eqref{eq:scaling} are then summed
as in~\eqref{eq:ensemble}. Section~\ref{sec:audit} explains why this replaces
the min max scaling of the source model. The two heads receive identical
treatment: the Adam optimiser with learning rate $5\times10^{-4}$ and weight
decay $10^{-5}$, minibatches of $64$, gradient clipping at norm $1$, up to
$150$ epochs with early stopping after $20$ epochs without improvement in the
training loss, and the learning rate halved after $8$ such epochs. Because
stopping is decided on the training loss, no data outside the training windows
influence any network. The weights used for prediction are an exponential
moving average with decay $\beta=0.999$, corrected for its initialisation,
\begin{equation}
\bar{\theta}_T=\frac{(1-\beta)\sum_{t=1}^{T}\beta^{\,T-t}\theta_t}{1-\beta^{T}},
\label{eq:ema}
\end{equation}
so that the average is taken over the weights actually visited and does not
depend on how long a component trained. Seeds are fixed per component. All
networks were trained on a single NVIDIA GeForce GTX~1650 graphics card.

\subsection{Evaluation measures}
Accuracy is reported as root mean squared error (RMSE), mean absolute error
(MAE), mean absolute percentage error (MAPE) and the coefficient of
determination $R^2$, following guidance on choosing error measures for
forecast evaluation~\cite{hewamalage2022}. Five benchmarks need no learning:
persistence, $\hat{x}_t=x_{t-1}$; a random walk with drift estimated on the
training data; moving averages of the last $5$ and $20$ values; and
exponential smoothing with span $10$. Two forecasts $a$ and $b$ are compared
with the MDM test~\cite{mdmindex2023}. With squared error loss differential
$d_t=e_{a,t}^2-e_{b,t}^2$ over $T$ test days and horizon $h=1$, the statistic is
\begin{equation}
\mathrm{MDM}=\sqrt{\frac{T+1-2h+h(h-1)/T}{T}}\;
\frac{\bar{d}}{\sqrt{\hat{\gamma}_0/T}},
\label{eq:mdm}
\end{equation}
where $\bar{d}$ is the mean and $\hat{\gamma}_0$ the variance of $d_t$, and
the statistic is referred to a Student $t$ distribution with $T-1$ degrees of
freedom. A positive value means that $a$ has the larger loss. The effect of the
protocol is summarised by the leakage inflation ratio
\begin{equation}
\rho=\frac{\mathrm{RMSE}_{\mathrm{B}}}{\mathrm{RMSE}_{\mathrm{A}}},
\label{eq:rho}
\end{equation}
computed on common test dates for each index, price series and head. A value
of $\rho=1$ means that decomposing the full series gave no advantage; a value
above one measures how much of the full series accuracy depends on future
data.

\subsection{Issues identified in the reproduced pipeline}\label{sec:audit}
The reproduction started from the source implementation and changed it only
where a change was needed for the comparison to be valid.
Tables~\ref{tab:audit_decomp} and~\ref{tab:audit_train} list every issue we
identified, its consequence, and how it was resolved, the first for the
decomposition and data handling stages and the second for training, so that a
reader can judge which of the differences between our numbers and the source's
are due to which change. Where the original behaviour was a modelling choice
rather than a fault, it was kept and reported.

\begin{table}[ht]
\caption{Issues identified in the decomposition and data handling stages of the source pipeline, and their resolution.}
\label{tab:audit_decomp}
\small
\begin{tabular}{@{}L{0.30\textwidth}L{0.39\textwidth}L{0.23\textwidth}@{}}
\toprule
Issue & Consequence & Resolution \\
\midrule
ICEEMDAN returned its components in reverse order &
VMD redecomposed the lowest frequency residual instead of the highest frequency mode &
Order corrected \\
Realisations with fewer modes were padded after the residual &
Averaging mixed one realisation's residual with another's mode &
Padding placed before the residual \\
Swarm leader compared against its own overwritten fitness &
The global best could not improve once it led &
Global best tracked separately \\
Walk forward component target copied from the last input value &
Component level scores trivially perfect; price level scores unaffected &
Target taken from the decomposition that includes day $t$ \\
Walk forward test days dropped when the mode count changed &
Up to $96.5\%$ of test days discarded (SSEC highest: $17$ of $479$ kept) &
Mode count reconciled; all test days kept \\
Walk forward histories of odd length trimmed at their newest end &
On alternate test days the input window ended at $t-2$, turning half the walk forward forecasts into two step forecasts &
Oldest observation left out instead, Eq.~\eqref{eq:walkforward}; every window ends at $t-1$ \\
Highest and lowest prices cleaned of missing rows separately &
The two series could refer to different trading days &
Missing rows removed jointly \\
Protocols split one day apart &
Row $i$ of each protocol refers to a different date &
Comparisons on common dates \\
Swarm objective decreases monotonically in $K$ and $1/\alpha$ &
Search returns the bounds $K=10$, $\alpha=100$ on every run &
Kept as published; reported \\
ICEEMDAN noise fixed at $0.2$ price units &
Noise is $2$ to $8\times10^{-4}$ of the series standard deviation &
Kept as published; reported \\
\botrule
\end{tabular}
\end{table}

\begin{table}[ht]
\caption{Issues identified in the training stage of the source pipeline, and their resolution.}
\label{tab:audit_train}
\small
\begin{tabular}{@{}L{0.30\textwidth}L{0.39\textwidth}L{0.23\textwidth}@{}}
\toprule
Issue & Consequence & Resolution \\
\midrule
Moving average of the weights initialised at the random initialisation &
After $930$ updates, $39\%$ of the averaged weights were still the random initialisation &
Bias correction, Eq.~\eqref{eq:ema} \\
Heads trained differently (TCN: $50$ epochs at learning rate $10^{-3}$, no early stopping, clipping or averaging; TRM: $150$ epochs at $5\times10^{-4}$ with all three, and a bias correction applied to it alone) &
Head comparison confounded with training treatment &
Identical training for both heads \\
Components min max scaled on their training rows &
The S\&P~500 trend component lay entirely above its training maximum throughout the test period; full series S\&P~500 high forecasts reached RMSE $198$ (TCN) and $915$ (TRM) against $47$ for persistence &
Windows taken relative to their last value, Eq.~\eqref{eq:scaling} \\
TRM head under min max scaling &
Constant forecasts for $11$ of the $16$ S\&P~500 high components &
Resolved by the same change \\
\botrule
\end{tabular}
\end{table}

\subsection{Reproducibility}
A single run of the study consists of five steps, and the counts below are
read from the files each step saves.
\begin{enumerate}
\item Download the unadjusted daily highest and lowest prices of both indices
once, remove rows with a missing value jointly, and cache the four series
($2513$ observations for each S\&P~500 series and $2425$ for each SSEC series).
\item For protocol A, decompose each full series once with $50$ noise
realisations, search $K$ and $\alpha$ with PSO, and window the components,
giving $497$ test days per S\&P~500 series and $479$ per SSEC series.
\item For protocol B, decompose the training portion with $20$ realisations,
fix $K$, $\alpha$ and $C$, and decompose every history from the split onward:
$498$ decompositions per S\&P~500 series and $480$ per SSEC series, $1956$ in
all.
\item Train one network per component, protocol, series and head, each with a
fixed seed: $16$ or $17$ components, two protocols, four series and two heads,
$262$ networks in all, producing $16$ forecast files.
\item Score every forecast file against the price on the dates shared by both
protocols ($496$ per S\&P~500 series and $478$ per SSEC series), together with
the five benchmarks and the MDM test.
\end{enumerate}
Five secondary experiments support the main comparison: the input diagnostic,
which asks whether the inputs alone predict the next day's change before any
network is trained (Table~\ref{tab:diagnostic}, Fig.~\ref{fig:mechanism}); the
realisation count check, which asks whether $20$ realisations in protocol B
change the components relative to $50$ (Section~\ref{sec:data}); the boundary
amplitude comparison, which asks how variable freshly decomposed components are
at the newest end of the history (Section~\ref{sec:results}); the input
persistence check, which asks what a network would score by copying its own
inputs (Section~\ref{sec:discussion}); and the archived run under the source
model's min max scaling, whose outcome is recorded in Table~\ref{tab:audit_train}.
Every table and figure in this paper is generated by script from the saved
forecast files.

\section{Results}\label{sec:results}

Five subsections follow: the information present in the inputs before any
network is trained, the accuracy of the forecasts under the two protocols, the
ratio between them, the comparison with benchmarks, and the comparison between
the two heads.

\subsection{Information about the next day in the inputs}
Before any network is trained, the inputs themselves already show how the two
protocols differ. Table~\ref{tab:diagnostic} relates the high frequency group,
the sum of the $K$ VMD modes on the last day of each input window, to the price
change that the forecast has to predict. Under full series decomposition the
two are strongly and negatively correlated on every series, with $r$ between
$-0.58$ and $-0.53$ ($p<10^{-35}$), and a linear regression on three input
features, fitted on half of the test dates and scored on the other half,
removes between $51\%$ and $60\%$ of the squared error of persistence. Daily
index changes should not allow anything of the kind. Fig.~\ref{fig:mechanism}
shows why it happens: when the full series is decomposed, the high frequency
value at day $t-1$ is shaped by the price at day $t$, so a rise on the next day
leaves day $t-1$ looking like a local trough and pulls its high frequency value
down. Under walk forward decomposition the same correlation lies between
$0.00$ and $0.10$, and the regression does worse than persistence on all four
series, by $0.9\%$ to $3.0\%$. The one correlation that reaches nominal
significance, $r=0.10$ for the SSEC highest price ($p=0.03$), carries no value
out of sample. The walk forward high frequency group is also $2.5$ to $3.3$
times more variable on the last day of the history than the full series group
on the same dates, which is the boundary behaviour that sifting is known to
need care with~\cite{emdtutorial2023}. Those large end values are noise rather
than signal, and a forecaster that works from history alone has to live with
them.

\begin{table}[ht]
\caption{Information about the next day's price change in the model inputs, measured without any network. $r$ is the correlation between the high frequency group at day $t-1$ and the change from $t-1$ to $t$; skill is the out of sample reduction in squared error that a linear regression on three input features achieves over persistence.}
\label{tab:diagnostic}
\centering
\begin{tabular}{@{}llrrr@{}}
\toprule
Series & Protocol & $r$ & $p$ & Skill \\
\midrule
S\&P~500 high & full series & ${-}0.526$ & $<0.001$ & $0.552$ \\
 & walk forward & $0.056$ & 0.21 & ${-}0.014$ \\
S\&P~500 low & full series & ${-}0.552$ & $<0.001$ & $0.509$ \\
 & walk forward & $0.045$ & 0.32 & ${-}0.030$ \\
SSEC high & full series & ${-}0.580$ & $<0.001$ & $0.596$ \\
 & walk forward & $0.098$ & 0.03 & ${-}0.015$ \\
SSEC low & full series & ${-}0.579$ & $<0.001$ & $0.556$ \\
 & walk forward & $0.002$ & 0.97 & ${-}0.009$ \\
\botrule
\end{tabular}
\end{table}

\begin{figure}[ht]
\centering
\figorbox{\textwidth}{figs/fig_mechanism.png}
\caption{Next day change of the S\&P~500 highest price plotted against the
high frequency group on the last day of the input window, for the $496$ test
dates both protocols share. Left, full series decomposition; right, walk
forward decomposition. The figure shows the shape of the relation that
Table~\ref{tab:diagnostic} summarises with one coefficient per series.}
\label{fig:mechanism}
\end{figure}

\subsection{Forecast accuracy under the two protocols}
Table~\ref{tab:accuracy} reports the accuracy of both heads under both
protocols, with persistence alongside, on the $496$ (S\&P~500) and $478$
(SSEC) test dates that the protocols share. Under full series decomposition
the forecasts are very accurate by any of the usual measures: RMSE between
$5.0$ and $10.7$ index points, MAPE between $0.066\%$ and $0.134\%$, and $R^2$
between $0.9992$ and $0.9999$. These values are typical of what decomposition
ensemble studies report. Under walk forward decomposition the same networks,
trained with the same code, seeds and budgets, reach RMSE between $42.5$ and
$68.5$ points, MAPE between $0.72\%$ and $0.83\%$, and $R^2$ between $0.981$
and $0.991$. Persistence sits between the two, with RMSE from $36.2$ to $57.9$
points and MAPE from $0.52\%$ to $0.67\%$. Fig.~\ref{fig:trace} shows the
first hundred test days of two series, where the difference is visible without
any statistic: the full series forecast lies on top of the price and turns
when the price turns, which is only possible because its inputs were computed
from the price it is forecasting, while the walk forward forecast follows the
price with a delay, as persistence does, and with additional error of its own.

\begin{table}[ht]
\caption{Out of sample accuracy on the dates tested by both protocols. Persistence forecasts each price with the previous day's value.}
\label{tab:accuracy}
\centering
\begin{tabular}{@{}llrrrr@{}}
\toprule
Series & Forecast & RMSE & MAE & MAPE (\%) & $R^2$ \\
\midrule
S\&P~500 high & Persistence & 46.72 & 33.57 & 0.523 & 0.9949 \\
 & TCN, full series & 5.83 & 4.23 & 0.066 & 0.9999 \\
 & TCN, walk forward & 60.66 & 46.61 & 0.724 & 0.9913 \\
 & TRM, full series & 7.25 & 5.37 & 0.084 & 0.9999 \\
 & TRM, walk forward & 60.47 & 46.37 & 0.720 & 0.9914 \\
\midrule
S\&P~500 low & Persistence & 57.92 & 42.34 & 0.669 & 0.9923 \\
 & TCN, full series & 7.53 & 5.38 & 0.085 & 0.9999 \\
 & TCN, walk forward & 68.52 & 51.48 & 0.808 & 0.9892 \\
 & TRM, full series & 10.74 & 7.21 & 0.116 & 0.9997 \\
 & TRM, walk forward & 65.25 & 49.39 & 0.776 & 0.9902 \\
\midrule
SSEC high & Persistence & 36.21 & 22.14 & 0.605 & 0.9883 \\
 & TCN, full series & 7.13 & 3.66 & 0.102 & 0.9995 \\
 & TCN, walk forward & 45.60 & 30.40 & 0.829 & 0.9815 \\
 & TRM, full series & 9.29 & 4.74 & 0.132 & 0.9992 \\
 & TRM, walk forward & 45.30 & 30.05 & 0.820 & 0.9817 \\
\midrule
SSEC low & Persistence & 36.91 & 24.12 & 0.666 & 0.9879 \\
 & TCN, full series & 4.99 & 3.22 & 0.091 & 0.9998 \\
 & TCN, walk forward & 42.52 & 29.29 & 0.806 & 0.9839 \\
 & TRM, full series & 7.82 & 4.73 & 0.134 & 0.9995 \\
 & TRM, walk forward & 42.54 & 29.29 & 0.807 & 0.9839 \\
\botrule
\end{tabular}
\end{table}

\begin{figure}[ht]
\centering
\figorbox{\textwidth}{figs/fig_trace.png}
\caption{Price, persistence and the two TCN forecasts over the first $100$
shared test days of the S\&P~500 highest price (top) and the SSEC highest price
(bottom). The figure adds the timing of the forecasts, which the error
measures in Table~\ref{tab:accuracy} do not show.}
\label{fig:trace}
\end{figure}

\subsection{How much of the accuracy depends on the protocol}
Table~\ref{tab:leakage} and Fig.~\ref{fig:rho} give the leakage inflation
ratio of Eq.~\eqref{eq:rho} for each series and head. Restricting the
decomposition to the past multiplies the RMSE by a factor between $4.9$ and
$10.4$, with a geometric mean of $7.2$ over the eight cases. The MDM statistic
of Eq.~\eqref{eq:mdm} for the walk forward forecast against the full series
forecast lies between $+4.5$ and $+12.4$, so every one of the eight
differences is significant at any conventional level. The ratio is larger for
the TCN head than for the TRM head on every series ($6.4$ to $10.4$ against
$4.9$ to $8.3$), because the TCN head achieved the lower full series error
while the two heads reach almost the same walk forward error.

\begin{table}[ht]
\caption{Effect of the decomposition protocol. $\rho$ is the walk forward RMSE (protocol B) divided by the full series RMSE (protocol A); a positive MDM statistic means the walk forward forecast has the larger squared error.}
\label{tab:leakage}
\centering
\begin{tabular}{@{}llrrrrr@{}}
\toprule
Series & Head & RMSE$_{\mathrm{A}}$ & RMSE$_{\mathrm{B}}$ & $\rho$ & MDM & $p$ \\
\midrule
S\&P~500 high & TCN & 5.83 & 60.66 & 10.40 & ${+}12.33$ & $<0.001$ \\
 & TRM & 7.25 & 60.47 & 8.34 & ${+}12.42$ & $<0.001$ \\
S\&P~500 low & TCN & 7.53 & 68.52 & 9.10 & ${+}10.08$ & $<0.001$ \\
 & TRM & 10.74 & 65.25 & 6.07 & ${+}10.22$ & $<0.001$ \\
SSEC high & TCN & 7.13 & 45.60 & 6.40 & ${+}4.81$ & $<0.001$ \\
 & TRM & 9.29 & 45.30 & 4.88 & ${+}4.51$ & $<0.001$ \\
SSEC low & TCN & 4.99 & 42.52 & 8.52 & ${+}6.24$ & $<0.001$ \\
 & TRM & 7.82 & 42.54 & 5.44 & ${+}5.96$ & $<0.001$ \\
\botrule
\end{tabular}
\end{table}

\begin{figure}[ht]
\centering
\figorbox{0.8\textwidth}{figs/fig_rho.png}
\caption{Leakage inflation ratio $\rho$ by series and head. A star marks a
significant MDM test of the walk forward forecast against the full series
forecast ($p<0.05$); the line at $\rho=1$ is where decomposing the full series
would give no advantage. The figure makes the consistency of the effect across
series and heads visible at a glance, which the rows of Table~\ref{tab:leakage}
do not.}
\label{fig:rho}
\end{figure}

\subsection{Comparison with benchmarks}
Table~\ref{tab:baselines} gives the RMSE of the five benchmarks. Persistence
and the drift forecast are indistinguishable, within $0.3\%$ of each other on
every series, and both are far ahead of any smoothing benchmark: the five day
moving average is $1.5$ to $1.6$ times worse, exponential smoothing $1.8$ to
$1.9$ times, and the twenty day moving average $2.7$ to $2.9$ times. On daily
index levels, the last observed price is the benchmark to beat.
Table~\ref{tab:dm} tests each model against it. Under full series
decomposition both heads beat persistence on all four series, with MDM
statistics between $-4.5$ and $-10.1$ and $p<0.001$ in every case. Under walk
forward decomposition both heads lose to persistence on all four series, with
MDM statistics between $+3.5$ and $+7.4$, again $p<0.001$ in every case. In
terms of RMSE, the causal forecasts are between $13\%$ and $30\%$ worse than
persistence. Of the eight model and series combinations, all eight beat
persistence when the decomposition sees the future and none does when it does
not.

\begin{table}[ht]
\caption{RMSE of the benchmark forecasts on the common test dates.}
\label{tab:baselines}
\centering
\begin{tabular}{@{}lrrrr@{}}
\toprule
Benchmark & S\&P~500 high & S\&P~500 low & SSEC high & SSEC low \\
\midrule
Persistence & 46.72 & 57.92 & 36.21 & 36.91 \\
Drift & 46.60 & 57.82 & 36.23 & 36.92 \\
MA-5 & 72.73 & 92.48 & 59.34 & 56.41 \\
MA-20 & 136.96 & 156.83 & 103.61 & 99.33 \\
EWMA-10 & 87.38 & 104.75 & 68.46 & 65.23 \\
\botrule
\end{tabular}
\end{table}

\begin{table}[ht]
\caption{Each model against persistence. A negative MDM statistic means the model has the smaller squared error.}
\label{tab:dm}
\centering
\begin{tabular}{@{}llrr@{}}
\toprule
Series & Forecast & MDM & $p$ \\
\midrule
S\&P~500 high & TCN, full series & ${-}9.94$ & $<0.001$ \\
 & TCN, walk forward & ${+}7.21$ & $<0.001$ \\
 & TRM, full series & ${-}9.88$ & $<0.001$ \\
 & TRM, walk forward & ${+}7.44$ & $<0.001$ \\
S\&P~500 low & TCN, full series & ${-}10.00$ & $<0.001$ \\
 & TCN, walk forward & ${+}5.03$ & $<0.001$ \\
 & TRM, full series & ${-}10.05$ & $<0.001$ \\
 & TRM, walk forward & ${+}4.45$ & $<0.001$ \\
SSEC high & TCN, full series & ${-}4.48$ & $<0.001$ \\
 & TCN, walk forward & ${+}4.26$ & $<0.001$ \\
 & TRM, full series & ${-}4.53$ & $<0.001$ \\
 & TRM, walk forward & ${+}3.54$ & $<0.001$ \\
SSEC low & TCN, full series & ${-}6.08$ & $<0.001$ \\
 & TCN, walk forward & ${+}4.10$ & $<0.001$ \\
 & TRM, full series & ${-}6.06$ & $<0.001$ \\
 & TRM, walk forward & ${+}3.60$ & $<0.001$ \\
\botrule
\end{tabular}
\end{table}

\subsection{TCN and TRM heads}
Table~\ref{tab:heads} compares the two heads under each protocol, using the
MDM statistic of the TCN forecast against the TRM forecast. Under full series
decomposition the TCN head is significantly better on all four series (MDM
between $-2.9$ and $-6.1$, $p<0.01$), which is the ranking the source
comparison reported. Under walk forward decomposition the picture changes: on
three series the difference is not significant ($p=0.56$, $0.52$ and $0.97$),
and on the fourth, the S\&P~500 lowest price, the TRM head is significantly
better (MDM $+4.2$, $p<0.001$). The two heads reach walk forward RMSE values
that differ by less than $0.7\%$ on three of the four series.

\begin{table}[ht]
\caption{TCN head against TRM head. A positive MDM statistic means the TCN head has the larger squared error.}
\label{tab:heads}
\centering
\begin{tabular}{@{}lrrrr@{}}
\toprule
Series & MDM, full series & $p$ & MDM, walk forward & $p$ \\
\midrule
S\&P~500 high & ${-}6.14$ & $<0.001$ & ${+}0.58$ & 0.559 \\
S\&P~500 low & ${-}4.61$ & $<0.001$ & ${+}4.22$ & $<0.001$ \\
SSEC high & ${-}2.94$ & 0.003 & ${+}0.65$ & 0.519 \\
SSEC low & ${-}4.07$ & $<0.001$ & ${-}0.04$ & 0.968 \\
\botrule
\end{tabular}
\end{table}

\section{Discussion}\label{sec:discussion}

The discussion has five parts: what the full series numbers measure, why the
causal forecasts trail persistence, what the head comparison shows, the
practices we recommend, and the limits of the evidence.

\subsection{What the full series numbers measure}
The full series results in Table~\ref{tab:accuracy} are of the same order as
those reported for this family of models: percentage errors around one tenth
of one percent and coefficients of determination above $0.999$. They were
produced by a pipeline whose implementation we audited line by line, with no
test data in any scaler, training set or stopping rule. They are nonetheless
not forecasts. Table~\ref{tab:diagnostic} shows that the inputs alone, with no
network, explain half of the next day's price change under this protocol, and
Table~\ref{tab:leakage} shows that the whole advantage over persistence
disappears when the decomposition is restricted to the past. The number that a
full series protocol measures is how efficiently a network can read the future
out of components that already contain it.

\subsection{Why the causal forecasts trail persistence}
The walk forward forecasts were not merely no better than persistence; they
were significantly worse on every series, by $13\%$ to $30\%$ in RMSE. Two
measurements locate the cause. First, the components at the newest end of a
freshly decomposed history are $2.5$ to $3.3$ times more variable than the same
components inside a full decomposition, which is the boundary behaviour of
sifting~\cite{emdtutorial2023}. Second, a forecaster that simply copied the
last value of each component window and summed the copies would have matched
persistence to within $0.7\%$ under either protocol, because the components
reconstruct the price. The trained networks therefore lost between $12\%$ and
$29\%$ relative to copying their own inputs. They were trained on windows from
the interior of the training decomposition, where component dynamics are
smooth, and applied to windows that end at a boundary, where they are not. This
is a mismatch between training and test inputs that the walk forward protocol
exposes rather than creates, and it applies equally to the source model and to
the many published variants that decompose the training portion once. Closing
it would require training on windows built the same way as the test windows,
at the cost of one extra decomposition per training day, roughly two thousand
per series, or an explicit boundary treatment inside the decomposition.

\subsection{The head comparison}
Under full series decomposition the TCN head was significantly more accurate
than the TRM head on all four series, which is the conclusion the source
comparison reached. Under walk forward decomposition the difference vanished
on three series and reversed on the fourth. The ranking of the two heads is
therefore a property of the evaluation protocol, not of the heads: when the
inputs contain the answer, the architecture that reads it out more efficiently
wins, and when they do not, the two are equivalent. This is consistent with the
observation that simple models hold their own against elaborate architectures
when forecasts are evaluated carefully~\cite{zeng2023,hewamalage2022}.

\subsection{Recommendations}
Three practices would have made the source result, and results like it,
interpretable.
\begin{enumerate}
\item Decompose only data that would have been available at the forecast
origin, and say so. Table~\ref{tab:diagnostic} gives a check that takes a few
seconds and needs no model.
\item Report persistence and the MDM statistic against it, since on daily
index levels it is the benchmark that matters.
\item Report the ratio $\rho$ of causal to noncausal error, so that readers
can see how much of a headline number depends on the protocol.
\end{enumerate}

\subsection{Limitations}
Five choices bound these conclusions.
\begin{enumerate}
\item The noise amplitude and the swarm objective were kept from the source
model, so the decomposition studied is ICEEMDAN at a small, fixed noise level
with $K$ and $\alpha$ at the bounds of their search ranges. Other settings
could change the size of the effect, although not the fact that a
decomposition of the full series has seen the prices it is asked to forecast.
\item The forecast horizon is one day, and the evidence covers two indices
over one ten year window with one seed per network.
\item Early stopping monitors the training loss, which keeps every network
away from test data but may leave individual components over or under
trained.
\item Input windows are expressed relative to their last value rather than
min max scaled, a change the S\&P~500 required but one that departs from the
source model, so our absolute errors are not directly comparable with the
errors reported there.
\item Walk forward training windows come from the interior of the training
decomposition while walk forward test windows end at a boundary, as discussed
above. The causal results therefore describe the protocol as it is commonly
run, and whether a boundary matched training scheme would narrow the gap to
persistence is untested here.
\end{enumerate}

\section{Conclusion}\label{sec:conclusion}
We reproduced a published decomposition ensemble stock forecaster, corrected
twelve implementation issues along the way, and evaluated it on the daily
highest and lowest prices of the S\&P~500 and the SSEC under two protocols that
differ only in whether the decomposition may see the future. With the series
decomposed once in full, both network heads reached the accuracy the
literature reports, MAPE near $0.1\%$ and $R^2$ above $0.999$, and beat
persistence on every series with $p<0.001$; a linear regression on the raw
inputs already explained half of the next day's price change, which is the
signature of information from day $t$ inside the features for day $t-1$. With
the decomposition restricted to the available history, the error rose by a
factor of $4.9$ to $10.4$, every forecast fell significantly behind
persistence, and the reported advantage of the convolutional head over the
recursive head disappeared. Decomposition ensemble results for daily index
levels should therefore be read as claims about the evaluation protocol until
they are shown to survive a causal decomposition and a persistence benchmark,
and the two checks proposed here, the input diagnostic and the leakage ratio,
make that demonstration inexpensive.

\backmatter

\begin{thebibliography}{99}

\bibitem{kumbure2022} M. M. Kumbure, C. Lohrmann, P. Luukka and J. Porras, ``Machine learning techniques and
data for stock market forecasting: A literature review,'' Expert Systems with
Applications, vol. 197, Art. no. 116659, 2022.

\bibitem{akimaemd2023} M. Ali, D. M. Khan, H. M. Alshanbari and A. A.-A. H. El-Bagoury, ``Prediction of
Complex Stock Market Data Using an Improved Hybrid EMD-LSTM Model,'' Applied Sciences,
vol. 13, no. 3, Art. no. 1429, 2023.

\bibitem{gruceemdan2023} C. Qi, J. Ren and J. Su, ``GRU Neural Network Based on CEEMDAN–Wavelet for Stock Price
Prediction,'' Applied Sciences, vol. 13, no. 12, Art. no. 7104, 2023.

\bibitem{ceemdanstock2025} J. K. Mutinda and A. Geletu, ``Stock Market Index Prediction Using
CEEMDAN‐LSTM‐BPNN‐Decomposition Ensemble Model,'' Journal of Applied Mathematics, vol.
2025, no. 1, Art. no. 7706431, 2025.

\bibitem{jiang2025} Q. Jiang, Z. Lin, J. Hu and X. Liu, ``Application of the VMD-CNN-BiLSTM-Attention
Model in Daily Price Forecasting of NYMEX Natural Gas Futures,'' Applied Sciences,
vol. 15, no. 20, Art. no. 11169, 2025.

\bibitem{wangwang2026} J.-W. Wang and X.-R. Wang, ``Intelligent Stock Price Forecasting Integrating Adaptive
VMD and Feature Selection with BiLSTM,'' Computational Economics, vol. 68, no. 4, pp.
3119-3153, 2026.

\bibitem{gong2024} H. Gong and H. Xing, ``Predicting the highest and lowest stock price indices: A
combined BiLSTM-SAM-TCN deep learning model based on re-decomposition,'' Applied Soft
Computing, vol. 167, Art. no. 112393, 2024.

\bibitem{zhao2025} Q. Zhao, H. Li, X. Liu and Y. Wang, ``A Hybrid Model of Multi-Head Attention Enhanced
BiLSTM, ARIMA, and XGBoost for Stock Price Forecasting Based on Wavelet Denoising,''
Mathematics, vol. 13, no. 16, Art. no. 2622, 2025.

\bibitem{ning2025} Z. Ning, G. Chen, J. Wang and W. Hu, ``Multi-Scale Decomposition and Hybrid Deep
Learning CEEMDAN-VMD-CNN-BiLSTM Approach for Wind Power Forecasting,'' Processes, vol.
13, no. 7, Art. no. 2046, 2025.

\bibitem{kapoor2023} S. Kapoor and A. Narayanan, ``Leakage and the reproducibility crisis in
machine-learning-based science,'' Patterns, vol. 4, no. 9, Art. no. 100804, 2023.

\bibitem{chen2022} Y. Chen, S. Yu, S. Islam, C. P. Lim and S. M. Muyeen, ``Decomposition-based wind power
forecasting models and their boundary issue: An in-depth review and comprehensive
discussion on potential solutions,'' Energy Reports, vol. 8, pp. 8805-8820, 2022.

\bibitem{emdleak2024} X. Yang, J. Li and X. Jiang, ``Research on information leakage in time series
prediction based on empirical mode decomposition,'' Scientific Reports, vol. 14, no.
1, Art. no. 28362, 2024.

\bibitem{wang2024} J. Wang, W. Yuan, S. Zhang, S. Cheng and L. Han, ``Implementing ultra-short-term wind
power forecasting without information leakage through cascade decomposition and
attention mechanism,'' Energy, vol. 312, Art. no. 133513, 2024.

\bibitem{yao2025} S. Yao, H. Zhu, X. Zhou, T. Peng and J. Zhang, ``LSTM Model Combined with Rolling
Empirical Mode Decomposition and Sample Entropy Reconstruction for Short-Term Wind
Speed Forecasting,'' Processes, vol. 13, no. 3, Art. no. 819, 2025.

\bibitem{pang2026} J. Pang and S. Dong, ``Stop unrealistic data preprocessing in wind speed forecasting:
approaches and discussions on preventing future data leakage,'' Applied Mathematical
Modelling, vol. 150, Art. no. 116376, 2026.

\bibitem{glasserman2024} P. Glasserman and C. Lin, ``Assessing Look-Ahead Bias in Stock Return Predictions
Generated by GPT Sentiment Analysis,'' The Journal of Financial Data Science, vol. 6,
no. 1, pp. 25-42, 2024.

\bibitem{pernagallo2025} G. Pernagallo, ``Random walks, Hurst exponent, and market efficiency,'' Quality \&
Quantity, vol. 59, no. S2, pp. 1097-1119, 2025.

\bibitem{finrel2024} H. H. Htun, M. Biehl and N. Petkov, ``Forecasting relative returns for S\&P 500 stocks
using machine learning,'' Financial Innovation, vol. 10, no. 1, Art. no. 118, 2024.

\bibitem{jolicoeur2025} A. Jolicoeur-Martineau, ``Less is More: Recursive Reasoning with Tiny Networks,''
arXiv preprint arXiv:2510.04871, 2025.

\bibitem{hewamalage2022} H. Hewamalage, K. Ackermann and C. Bergmeir, ``Forecast evaluation for data
scientists: common pitfalls and best practices,'' Data Mining and Knowledge Discovery,
vol. 37, no. 2, pp. 788-832, 2023.

\bibitem{mdmindex2023} X. Xu and Y. Zhang, ``China mainland new energy index price forecasting with the
neural network,'' Energy Nexus, vol. 10, Art. no. 100210, 2023.

\bibitem{emdtutorial2023} C. van Jaarsveldt, G. W. Peters, M. Ames and M. Chantler, ``Tutorial on Empirical Mode
Decomposition: Basis Decomposition and Frequency Adaptive Graduation in Non-Stationary
Time Series,'' IEEE Access, vol. 11, pp. 94442-94478, 2023.

\bibitem{iceemdanbridge2022} R. M. Delgadillo and J. R. Casas, ``Bridge damage detection via improved completed
ensemble empirical mode decomposition with adaptive noise and machine learning
algorithms,'' Structural Control and Health Monitoring, vol. 29, no. 8, 2022.

\bibitem{ceemdanpaa2022} L. Hu, L. Wang, Y. Chen, N. Hu and Y. Jiang, ``Bearing Fault Diagnosis Using Piecewise
Aggregate Approximation and Complete Ensemble Empirical Mode Decomposition with
Adaptive Noise,'' Sensors, vol. 22, no. 17, Art. no. 6599, 2022.

\bibitem{vmdmssr2022} Z. Meng, X. Wang, J. Liu and F. Fan, ``An adaptive spectrum segmentation-based
optimized VMD method and its application in rolling bearing fault diagnosis,''
Measurement Science and Technology, vol. 33, no. 12, Art. no. 125107, 2022.

\bibitem{psovmdtcn2023} G. Geng, Y. He, J. Zhang, T. Qin and B. Yang, ``Short-Term Power Load Forecasting
Based on PSO-Optimized VMD-TCN-Attention Mechanism,'' Energies, vol. 16, no. 12, Art.
no. 4616, 2023.

\bibitem{damasevicius2024} R. Dama\v{s}evi\v{c}ius, L. Jovanovic, A. Petrovic, M. Zivkovic, N. Bacanin, D.
Jovanovic and M. Antonijevic, ``Decomposition aided attention-based recurrent neural
networks for multistep ahead time-series forecasting of renewable power generation,''
PeerJ Computer Science, vol. 10, Art. no. e1795, 2024.

\bibitem{kreuzer2025} T. Kreuzer, J. Zdravkovic and P. Papapetrou, ``Unpacking the trend: decomposition as a
catalyst to enhance time series forecasting models,'' Data Mining and Knowledge
Discovery, vol. 39, no. 5, Art. no. 54, 2025.

\bibitem{zeng2023} A. Zeng, M. Chen, L. Zhang and Q. Xu, ``Are Transformers Effective for Time Series
Forecasting?'' Proceedings of the AAAI Conference on Artificial Intelligence, vol. 37,
no. 9, pp. 11121-11128, 2023.

\bibitem{makridakis2022} S. Makridakis, E. Spiliotis, V. Assimakopoulos, A.-A. Semenoglou, G. Mulder and K.
Nikolopoulos, ``Statistical, machine learning and deep learning forecasting methods:
Comparisons and ways forward,'' Journal of the Operational Research Society, vol. 74,
no. 3, pp. 840-859, 2023.

\bibitem{olorunnimbe2022} K. Olorunnimbe and H. Viktor, ``Deep learning in the stock market—a systematic survey
of practice, backtesting, and applications,'' Artificial Intelligence Review, vol. 56,
no. 3, pp. 2057-2109, 2023.

\bibitem{rnnreview2024} I. D. Mienye, T. G. Swart and G. Obaido, ``Recurrent Neural Networks: A Comprehensive
Review of Architectures, Variants, and Applications,'' Information, vol. 15, no. 9,
Art. no. 517, 2024.

\bibitem{colominas2014} M. A. Colominas, G. Schlotthauer and M. E. Torres, ``Improved complete ensemble EMD: A
suitable tool for biomedical signal processing,'' Biomedical Signal Processing and
Control, vol. 14, pp. 19-29, 2014.

\bibitem{dragomiretskiy2014} K. Dragomiretskiy and D. Zosso, ``Variational Mode Decomposition,'' IEEE Transactions
on Signal Processing, vol. 62, no. 3, pp. 531-544, 2014.

\bibitem{psosurvey2022} T. M. Shami, A. A. El-Saleh, M. Alswaitti, Q. Al-Tashi, M. A. Summakieh and S.
Mirjalili, ``Particle Swarm Optimization: A Comprehensive Survey,'' IEEE Access, vol.
10, pp. 10031-10061, 2022.

\end{thebibliography}

\bmhead{Supplementary information} Not applicable.

\bmhead{Acknowledgements} Not applicable.

\bmhead{Data availability} The study uses the public daily highest and lowest prices of
the S\&P~500 (ticker \texttt{\textasciicircum GSPC}) and the Shanghai Stock Exchange
Composite Index (\texttt{000001.SS}), retrieved once from Yahoo Finance and cached so
that every experiment uses the same arrays. The cached series, the decompositions under
both protocols, every forecast file and the scripts that generate every table and figure
in this paper are available at
\url{https://github.com/prayag-1771/Decomposition-Ensemble}.

\section*{Declarations}

\begin{itemize}
\item \textbf{Funding} \; No funding was received for this research.
\item \textbf{Conflict of interest/Competing interests} \; The authors declare that
they have no competing interests.
\item \textbf{Ethics approval} \; Not applicable. The study analyses publicly available
stock index prices only and involves no human participants, human data or animals.
\item \textbf{Consent to participate} \; Not applicable.
\item \textbf{Consent for publication} \; All authors consent to the publication of this
manuscript.
\item \textbf{Availability of data and materials} \; The cached price series, the
decomposed datasets under both protocols and all forecast files are available at
\url{https://github.com/prayag-1771/Decomposition-Ensemble}.
\item \textbf{Code availability} \; The code for the two decomposition protocols, the
training of the component networks, the benchmark forecasts and statistical tests, and
the scripts that generate every table and figure from the saved result files is
available at \url{https://github.com/prayag-1771/Decomposition-Ensemble}.
\item \textbf{Authors' contributions} \; Prayag Sharma implemented the training
pipeline and the input scaling, trained all component networks, carried out the audit
of the reproduced pipeline, performed the leakage diagnostic and the analysis of the
results, produced the tables and figures, and drafted the manuscript. Veer Vikram
Saxena acquired and cleaned the price series, implemented the two stage decomposition
and both decomposition protocols, and produced and validated the decomposed datasets.
Ayush Srivastava implemented the benchmark forecasts, the error measures and the
modified Diebold and Mariano test, and carried out the evaluation of the forecasts
against the benchmarks. Ravi Tiwari supervised the work, guided the design of the study
and reviewed and edited the manuscript. All authors read and approved the final
manuscript.
\end{itemize}

\end{document}
